\textbf{Phần 1. Trắc nghiệm lựa chọn}
\setcounterpageref{ex}{0}

\begin{ex} %[3X16N80]
Tung một con xúc xắc 6 mặt, khả năng xuất hiện mặt 8 chấm là
\choice
{Chắc chắn xảy ra}
{Có thể xảy ra}
{Rất dễ xảy ra}
{\True Không thể xảy ra}
\loigiai{
Xúc xắc chỉ có từ 1 đến 6 chấm nên không thể có 8 chấm.\\
\textbf{Đáp án của câu hỏi là D}
}
\end{ex}

\begin{ex} %[3X16N80]
Cho bảng số liệu số học sinh yêu thích môn học: Toán: 18 bạn, Tiếng Việt: 15 bạn, Tiếng Anh: 12 bạn. Môn học được nhiều bạn thích nhất là
\choice
{\True Toán}
{Tiếng Việt}
{Tiếng Anh}
{Ba môn bằng nhau}
\loigiai{
Môn Toán có 18 bạn thích (nhiều nhất).\\
\textbf{Đáp án của câu hỏi là A}
}
\end{ex}

\begin{ex} %[3X16H80]
Tổng số học sinh tham gia khảo sát môn học trên là
\choice
{$40$ bạn}
{$50$ bạn}
{$42$ bạn}
{\True $45$ bạn}
\loigiai{
$18 + 15 + 12 = 45$ bạn.\\
\textbf{Đáp án của câu hỏi là D}
}
\end{ex}

\begin{ex} %[3X16H80]
Số bạn thích môn Toán nhiều hơn Tiếng Anh là
\choice
{$5$ bạn}
{\True $6$ bạn}
{$7$ bạn}
{$8$ bạn}
\loigiai{
$18 - 12 = 6$ bạn.\\
\textbf{Đáp án của câu hỏi là B}
}
\end{ex}

\begin{ex} %[3X16H80]
Trong hộp có 5 viên bi xanh và 5 viên bi đỏ. Rút ngẫu nhiên 1 viên bi. Sự kiện \textit{"Rút được viên bi màu xanh"} là sự kiện
\choice
{Chắc chắn xảy ra}
{Không thể xảy ra}
{Luôn luôn xảy ra}
{\True Có thể xảy ra}
\loigiai{
Có thể rút được bi xanh hoặc bi đỏ.\\
\textbf{Đáp án của câu hỏi là D}
}
\end{ex}

\begin{ex} %[3X16H80]
Trong hộp chỉ có toàn bi vàng. Rút ngẫu nhiên 1 viên bi thì sự kiện nào \textbf{chắc chắn} xảy ra?
\choice
{Rút được bi xanh}
{\True Rút được bi vàng}
{Rút được bi đỏ}
{Rút được bi trắng}
\loigiai{
Chắc chắn rút được bi vàng.\\
\textbf{Đáp án của câu hỏi là B}
}
\end{ex}

\begin{ex} %[3X16V80]
Thống kê điểm bài kiểm tra cuối kì của tổ 2: Điểm 10: 4 bạn, Điểm 9: 5 bạn, Điểm 8: 3 bạn. Tổng số bạn đạt điểm từ 9 trở lên là
\choice
{$8$ bạn}
{$10$ bạn}
{\True $9$ bạn}
{$7$ bạn}
\loigiai{
$4 + 5 = 9$ bạn.\\
\textbf{Đáp án của câu hỏi là C}
}
\end{ex}

\begin{ex} %[3X16V80]
Tổ 2 có tất cả số học sinh là
\choice
{\True $12$ bạn}
{$11$ bạn}
{$13$ bạn}
{$10$ bạn}
\loigiai{
$4 + 5 + 3 = 12$ bạn.\\
\textbf{Đáp án của câu hỏi là A}
}
\end{ex}

\begin{ex} %[3X16V80]
Tung đồng xu 25 lần, mặt sấp xuất hiện 14 lần. Số lần mặt ngửa xuất hiện là
\choice
{$10$ lần}
{$12$ lần}
{\True $11$ lần}
{$13$ lần}
\loigiai{
$25 - 14 = 11$ lần.\\
\textbf{Đáp án của câu hỏi là C}
}
\end{ex}

\begin{ex} %[3X16V80]
Trong hộp có 10 chiếc thẻ từ 1 đến 10. Rút 1 thẻ. Sự kiện \textit{"Rút được thẻ ghi số có 1 chữ số"} là
\choice
{Chắc chắn xảy ra}
{\True Có thể xảy ra}
{Không thể xảy ra}
{Luôn luôn xảy ra}
\loigiai{
Có thể rút từ 1 đến 9 (1 chữ số) hoặc số 10 (2 chữ số).\\
\textbf{Đáp án của câu hỏi là B}
}
\end{ex}

\begin{ex} %[3X16V80]
Sự kiện \textit{"Rút được thẻ ghi số 15"} là
\choice
{\True Không thể xảy ra}
{Chắc chắn xảy ra}
{Có thể xảy ra}
{Rất dễ xảy ra}
\loigiai{
Thẻ chỉ có từ 1 đến 10 nên không thể rút được số 15.\\
\textbf{Đáp án của câu hỏi là A}
}
\end{ex}

\begin{ex} %[3X16C80]
Một hộp có 8 viên kẹo chanh và 6 viên kẹo cam. Cần lấy ra ít nhất bao nhiêu viên kẹo (không nhìn) để chắc chắn lấy được ít nhất 1 viên kẹo cam?
\choice
{$8$ viên}
{$10$ viên}
{\True $9$ viên}
{$7$ viên}
\loigiai{
Lấy hết 8 kẹo chanh + 1 kẹo nữa $= 9$ viên.\\
\textbf{Đáp án của câu hỏi là C}
}
\end{ex}

\textbf{Phần 2. Trắc nghiệm đúng sai}
\setcounterpageref{ex}{0}

\begin{ex} %[3X16H80]
Bảng số liệu số quả bóng trong kho: Bóng đá: 40 quả, Bóng rổ: 25 quả, Bóng chuyền: 35 quả. Xét tính đúng sai:
\choiceTF
{\True Bóng đá có số lượng nhiều nhất (40 quả).}
{\True Tổng số bóng trong kho là 100 quả ($40 + 25 + 35 = 100$).}
{Bóng chuyền ít hơn bóng rổ 10 quả.}
{\True Bóng chuyền nhiều hơn bóng rổ 10 quả ($35 - 25 = 10$).}
\loigiai{
\begin{itemize}
    \item Bóng đá nhiều nhất: Đúng.
    \item Tổng: $40 + 25 + 35 = 100$ quả. Đúng.
    \item Bóng chuyền (35) nhiều hơn bóng rổ (25). Ý 3 Sai, Ý 4 Đúng.
\end{itemize}
\begin{itemize}
    \item a) \textbf{Đúng}.
    \item b) \textbf{Đúng}.
    \item c) \textbf{Sai}.
    \item d) \textbf{Đúng}.
\end{itemize}
}
\end{ex}

\begin{ex} %[3X16V80]
Gieo một con xúc xắc 6 mặt. Xét tính đúng sai:
\choiceTF
{\True Có thể xuất hiện mặt 3 chấm.}
{\True Chắc chắn số chấm xuất hiện nằm từ 1 đến 6.}
{Chắc chắn xuất hiện mặt 6 chấm.}
{\True Không thể xuất hiện mặt 0 chấm.}
\loigiai{
\begin{itemize}
    \item Có thể mặt 3: Đúng.
    \item Chắc chắn từ 1 đến 6: Đúng.
    \item Không chắc chắn mặt 6. Sai.
    \item Không thể xuất hiện mặt 0 chấm: Đúng.
\end{itemize}
\begin{itemize}
    \item a) \textbf{Đúng}.
    \item b) \textbf{Đúng}.
    \item c) \textbf{Sai}.
    \item d) \textbf{Đúng}.
\end{itemize}
}
\end{ex}

\textbf{Phần 3. Trả lời ngắn}
\setcounterpageref{ex}{0}

\begin{ex} %[3X16H80]
Một nhóm thu gom được 3 nhóm (mỗi nhóm 5 bao rác) và 4 bao rác rời. Tính tổng số bao rác thu gom được.
\shortans{19}
\loigiai{
\textbf{Đáp số:} $19$\\
$3 \times 5 + 4 = 19$ bao.
}
\end{ex}

\begin{ex} %[3X16H80]
Tung đồng xu 50 lần, có 28 lần mặt sấp. Hỏi có bao nhiêu lần xuất hiện mặt ngửa?
\shortans{22}
\loigiai{
\textbf{Đáp số:} $22$\\
$50 - 28 = 22$ lần.
}
\end{ex}

\begin{ex} %[3X16V80]
Thống kê số sách đọc trong hè của 3 bạn: An đọc 12 cuốn, Bình đọc 15 cuốn, Cúc đọc 18 cuốn. Tính tổng số sách cả 3 bạn đã đọc.
\shortans{45}
\loigiai{
\textbf{Đáp số:} $45$\\
$12 + 15 + 18 = 45$ cuốn.
}
\end{ex}

\begin{ex} %[3X16V80]
Cúc đọc nhiều hơn An bao nhiêu cuốn sách?
\shortans{6}
\loigiai{
\textbf{Đáp số:} $6$\\
$18 - 12 = 6$ cuốn.
}
\end{ex}

\begin{ex} %[3X16V80]
Trong hộp có 12 viên bi đỏ và 1 viên bi xanh. Cần lấy ra ít nhất bao nhiêu viên bi để chắc chắn có 1 viên bi màu đỏ?
\shortans{2}
\loigiai{
\textbf{Đáp số:} $2$\\
$1 + 1 = 2$ viên bi.
}
\end{ex}

\begin{ex} %[3X16C80]
Một túi có 7 quả bóng xanh và 5 quả bóng đỏ. Cần lấy ra ít nhất bao nhiêu quả bóng để chắc chắn có ít nhất 1 quả bóng màu đỏ?
\shortans{8}
\loigiai{
\textbf{Đáp số:} $8$\\
$7 + 1 = 8$ quả bóng.
}
\end{ex}
